\section{The corrected fifth theorem, with all domains visible}

\begin{theorem}[C44-T05: synthesis of H9-T1, H9-T2, H9-T7,
H9-T17 and H9-T18]\label{thm:final-th5}
Use rational Chow groups and the signed geometric intersection height
of JAC1964 H9-T1. The following is the consolidated form of cote 44,
scan 004, Th.~5. The formulas below are unconditional under their
stated admissibility and geometric hypotheses. Each assertion that
uses the abelian Hodge standard conjecture is explicitly conditional
on the unrefereed 032/Milne inputs.

\begin{enumerate}
\item Let $f:\mathcal X\to C$ be projective with $\mathcal X$ smooth
of dimension $2m$, $C$ a smooth projective curve, and smooth generic
fibre $Y$. Let $x_0,y_0$ extend numerically trivial cycles on $Y$.
For each relevant fibre let $W_s$ be its vertical numerical cycle
space, $R_s$ its radical, and $M_s$ the nonsingular matrix on a chosen
basis of $W_s/R_s$. Suppose the extensions annihilate $R_s$, and
that these spaces account for all vertical pairings. Write $a_s,b_s$
for their pairing vectors. Then admissible corrections exist and
\[
 h_m(x_0|_Y,y_0|_Y)=(-1)^m
 \left(\deg(x_0y_0)-\sum_s a_s^{t}M_s^{-1}b_s\right).
 \tag{C44.T05.1}
\]
If only a smaller vertical space has been used, this expression is
the corrected intersection, and is a height only after the remaining
vertical orthogonality is verified. Failure to annihilate a vertical
radical is an obstruction to correction, not a term one may omit.

\item For a Lefschetz pencil on a smooth $2m$-fold $X$, if
$H^{2m}(Y_s)\to H^{2m}(Y_{\bar\eta_s})$ is injective at every
singular fibre, all corrections vanish for ambient cycles whose
generic restrictions are homologically trivial:
\[
                  h_m(j^*z,j^*z')=(-1)^m\int_Xzz'.
 \tag{C44.T05.2}
\]
The same injectivity makes any global extension of a homologically
trivial generic cycle admissible. In the ordinary Picard--Lefschetz
setting the condition is equivalent to nonzero vanishing cycles;
SGA 7 condition $(A)$ is the characteristic-independent formulation
used in the predecessor. For genuine abelian pencils it holds by
JAC-P07. Conditional on 032/Milne, (C44.T05.2) is positive definite
on primitive ambient algebraic cohomology.

\item If $X$ is abelian, let $B$ be its smooth axis and use
$\mathcal H,\mathcal H_0,A_X,A_B$ of H9-T7, (H9.17). Then
\[
 \mathcal H/\mathcal H_0\simeq A_X\mathbin\perp A_B,
 \qquad h_m\simeq
 \left((-1)^m\int_Xzz'\right)\mathbin\perp
 \left((-1)^{m-1}\int_Bww'\right).
 \tag{C44.T05.3}
\]
In particular $h_Bw=0$ for $w\in A_B$. Conditional on 032/Milne,
semipositivity on the whole $\mathcal H$ is equivalent to the weak
axis sign on $A_B$. If the axis form is positive definite, the
radical is exactly $\mathcal H_0$. If it is only semidefinite,
the quotient of the radical by $\mathcal H_0$ is the radical of
that axis form; equality with $\mathcal H_0$ is not inferred.
The axis sign is known for $m=2$, for characteristic zero, and in
the generic-axis vanishing cases H9-T8, H9-T8a and JAC-P09.
C44-A04 and C44-A06 give additional explicit positive subspaces,
not a sign theorem for every special axis.

\item For a smooth threefold generic fibre, on the admissible
homologically trivial domain let
\[
 \mathcal H_h^{\rm Alb}=\{z:\operatorname{alb}(hz)=0\},\qquad
 \psi_h(D)=\operatorname{alb}(h^2D),\qquad
 D_z=\psi_h^{-1}\operatorname{alb}(hz).
\]
The rational isogeny $\psi_h:\Pic^0(Y)\to\Alb(Y)$ gives the
orthogonal decomposition
\[
                    z=(z-hD_z)+hD_z.
 \tag{C44.T05.4}
\]
The first term lies in $\mathcal H_h^{\rm Alb}$; the second has
negative definite height modulo the constant trace. In
characteristic zero the first height is semidefinite by Zhang,
Theorem 1, with the precise admissible-model equality criterion
of H9-T17, (H9.77)--(H9.80). This sign statement does not extend
to the whole homologically trivial group after an arbitrary
finite cover: H9-T17 gives the exact negative value $-12d^3$ on
the displayed $E^4$ example, in every characteristic.

\item For an abelian-fourfold pencil and a finite separable cover
$C\to\mathbf P^1$ unramified over the critical values, let
$k:\mathcal X\hookrightarrow A\times C$ be the smooth incidence
fourfold. All homologically trivial generic cycles are admissible.
On $\mathcal H_h^{\rm Alb}$ the primitive correction of H9-T18
gives
\[
              h_2(z,z)=\int_{A\times C}a^2\lambda
                              +\int_{\mathcal X}v^2,
 \quad \lambda^2a=0,\quad k_*v=0.
 \tag{C44.T05.5}
\]
Conditional on 032/Milne the first term is positive definite on
its algebraic domain. The remaining assertion is precisely the
nonnegativity of the algebraic primitive form on $\ker k_*$.
C44-C01 splits off its positive old surface-axis part; the new
trace-kernel part is the remaining special-cover gate. The
generic vanishing results H9-T19--T26 and C44-C02--C04 give
the norm identities in their explicitly stated parameter scopes.
They do not establish this sign for every closed special cover.
\end{enumerate}
\end{theorem}

\begin{proof}
This is an assembly of the specified predecessor results, not a new
proof of them. Part (1) is H9-T1, (H9.3)--(H9.5), equal to
SIGN\_NEXT SN-44-01 after its normalization, with vertical blocks
separated by disjoint fibres. Part (2) is H9-T2, (H9.6)--(H9.7),
and its full-extension argument at the beginning of H9-T7;
the Picard--Lefschetz equivalence is SN-44-03. Part (3) is the
isometry H9-T7, (H9.18), with the elementary radical statement
of an orthogonal sum and the conditional ambient sign. Part (4)
is H9-T17, (H9.79)--(H9.84), including both of Zhang's hypotheses
and Kahn's trace quotient. Part (5) is H9-T12 admissibility and
H9-T18, (H9.86)--(H9.91), followed by the independently checked
old/new decomposition C44-C01. No step asserts algebraic
equivalence from homological triviality.
\end{proof}

\paragraph{The incidence higher Picard variety and the arithmetic reading.}
Grothendieck's $P^m$ in cote 44, scans 002 and 004, is attached
to algebraically trivial cycles via incidence. A homologically
trivial class in the displayed larger Chow groups is not thereby
a point of $P^m$. The Neron--Tate interpretation on a primitive
incidence quotient requires the actual regular homomorphism and
the height comparison; these are the separate inputs of
JAC-T4.1--T4.5 and SN-44-01. The full-middle-factor polarization
is refuted by JAC-P03, and the unrestricted covered height is
refuted by H9-T17. Neither counterexample refutes the primitive
axis sign asked here. The normalization and the Albanese kernel
implement the primitive-decomposition margin of scan 008.

In cote 16, scans 049--050, the algebraicity of Lefschetz operators
is condition $C=C_G=B$, while the height signs use $I$ on the
specified primitive algebraic spaces. Restoring the Lefschetz
factor in $A''$ (SN-C03--C04) is a separate issue from these
signs. Theorem C44-T05 does not deduce $I$ from $C$ or identify
cohomological semisimplicity with injectivity on rational Chow
groups. In the SGA framework, the local input is SGA 7 II,
XVIII 5.3.2, 5.3.4--5.3.5 and 6.2; the variable irreducibility
is XVIII 6.7.2. In the Chow-motive framework the global terms
are the blowup summand $h(B)(1)$ and the old/new finite-cover
projectors. These locators and their roles are retained from
the checked predecessor arguments.
